% GENERATED FILE. Edit canonical lessons, then run npm run generate:books.
% canonical-source-hash: fnv1a64:ac09d2ee6f6b2371
% canonical-lessons: TA-C285-iram, TA-S141-letter-ii, TA-C285-atan-piraku, TA-C285-vantu, TA-C285-mazhai-naal, TA-R285-a-story-in-order

\chapter{Telling a Story in Order}
\label{ch:ta-telling-a-story-in-order}
\hlchaptermodality{285}

\begin{chapteropening}
\textbf{By the end of this chapter:} \emph{I can tell a short story in order with mutalil, ataṉ piṟaku and kaṭaiciyil, join actions with vantu and pōy, read a connected story of about eighty words, and read the letter \ta{ஈ} in īram.}

Complete the last lesson of chapter 285: tell a story in order with first, after that, having done and in the end, and read and write \ta{ஈ}.
\end{chapteropening}

\section[\texorpdfstring{\ta{ஈரம்} (īram)}{ஈரம் (īram)}]{\ta{ஈரம்} (īram) — dampness, moisture}

\label{lesson:TA-C285-iram}



\begin{quote}
Before the new one: say \textquotedblleft{}It was raining.\textquotedblright{} (\emph{maḻai peytukoṇṭiruntatu}.) Then the words for a shirt and a cloth. (\emph{caṭṭai}, \emph{tuṇi}.)
\end{quote}

\subsection*{You'll want to know: \ta{ஈரம்}}
\textbf{\ta{ஈரம்}} — \emph{īram} — \textquotedblleft{}dampness, moisture\textquotedblright{}.

\begin{quote}
\textbf{\ta{மழை} \ta{பெய்தது}. \ta{என்} \ta{சட்டை} \ta{ஈரமாக} \ta{இருக்கிறது}.}
\end{quote}

\begin{quote}
\emph{maḻai peytatu. eṉ caṭṭai īramāka irukkiṟatu.} — \textquotedblleft{}It rained. My shirt is wet.\textquotedblright{}
\end{quote}

\textbf{\ta{ஈரம்}} names the wetness itself. To say a thing \textbf{is} wet, add the \textbf{-\ta{ஆக}} you know from \textbf{\ta{நன்றாக}} and \textbf{\ta{மெதுவாக}}, then \textbf{\ta{இருக்கிறது}}: \textbf{\ta{ஈரமாக} \ta{இருக்கிறது}}, \textquotedblleft{}it is damp\textquotedblright{}. In front of a noun it takes \textbf{-\ta{ஆன}}: \textbf{\ta{ஈரமான} \ta{துணி}} (\emph{īramāṉa tuṇi}, \textquotedblleft{}a wet cloth\textquotedblright{}).

Say the first sound long: \emph{ī-ram}. It is the long \emph{ī} you know as the sign \ta{ீ} on \textbf{\ta{நீ}}, here standing at the very front of the word. Standing alone, it is a letter of its own, one you have not met as a letter yet. The next lesson takes it out.

\begin{cousinweb}
\textbf{\ta{ஈரம்}} (\emph{īram}) is native Dravidian. DEDR 3690 lists Tamil \emph{īram}, \textquotedblleft{}wet, moisture\textquotedblright{}, in the same entry as \emph{nīr}, \textquotedblleft{}water\textquotedblright{}, the \textbf{\ta{நீர்}} at the end of \textbf{\ta{தேநீர்}}, and sets Malayalam \emph{īram} and Kannada \emph{īra}, \textquotedblleft{}moisture\textquotedblright{}, beside it.
\end{cousinweb}

\subsection*{Your turn}
Say these aloud:

\begin{itemize}
\raggedright
  \item \emph{īram}, slowly, then at speed
  \item \textquotedblleft{}It rained. My shirt is wet.\textquotedblright{} — \emph{maḻai peytatu. eṉ caṭṭai īramāka irukkiṟatu.}
  \item \textquotedblleft{}a wet cloth\textquotedblright{} — \emph{īramāṉa tuṇi}
\end{itemize}

\begin{tcolorbox}[breakable,colback=teal!4,colframe=teal!35!black,title={Before you move on}]
What does \textbf{\ta{ஈரம்}} mean? (\textbf{Dampness, moisture.}) \textquotedblleft{}It is wet\textquotedblright{}? (\textbf{\ta{ஈரமாக} \ta{இருக்கிறது}}.) \textquotedblleft{}A wet cloth\textquotedblright{}? (\textbf{\ta{ஈரமான} \ta{துணி}}.) Which word does DEDR list it with? (\textbf{\ta{நீர்}}, \textquotedblleft{}water\textquotedblright{}.)
\end{tcolorbox}
\section[\texorpdfstring{\ta{ஈ} (ī)}{ஈ (ī)}]{\ta{ஈ} (ī) — the long i, standing alone, taken from \ta{ஈரம்}}

\label{lesson:TA-S141-letter-ii}



\begin{quote}
Before the new one: read \textbf{\ta{இ}} aloud (\emph{i}), then \textbf{\ta{நீ}} (\emph{nī}), whose long \emph{ī} is the sign \ta{ீ}.

Now say \textbf{\ta{ஈரம்}} (\emph{īram}), \textquotedblleft{}dampness\textquotedblright{}. You met it a moment ago. Its first character is the letter of this lesson.
\end{quote}

\subsection*{Script you'll notice: \ta{ஈ}}
\textbf{\ta{ஈ}} — \emph{ī}.

It opens \textbf{\ta{ஈரம்}} (\emph{īram}). It is the long partner of \textbf{\ta{இ}} (\emph{i}), as \textbf{\ta{ஊ}} is of \textbf{\ta{உ}}, and \textbf{\ta{ஏ}} of \textbf{\ta{எ}}. After a consonant, you have read this vowel for a long time as the sign \ta{ீ}: \textbf{\ta{நீ}}, \textbf{\ta{தேநீர்}}. \textbf{\ta{ஈ}} is the same vowel when nothing comes before it, at the front of a word.

Its shape does not grow out of \textbf{\ta{இ}} the way \textbf{\ta{ஊ}} grows out of \textbf{\ta{உ}}. In print the two look nothing alike: \textbf{\ta{இ}} is all curves, a small loop inside a larger one; \textbf{\ta{ஈ}} is straight lines and dots. In the print of this book, \textbf{\ta{ஈ}} has a flat bar along the top and two straight legs hanging down from it, one from the left end and one from the middle. The right half of the bar has no leg. Two round dots sit halfway down, one between the two legs and one under the right half of the bar.

\subsection*{Writing: \ta{ஈ} — copy the printed shape}
Put your pen beside \textbf{\ta{ஈ}} and copy the shape you can see, slowly and large. Check the copy against the print: the bar, the left leg, the middle leg, the two dots.

\begin{quote}
This book does not yet tell you where to start \ta{ஈ} or which way the pen travels. That is real knowledge, taught with real variation from school to school, and it is not written down here until it can be written down with a source. Copying the printed shape needs no such source, and it is how the shape becomes yours in the meantime.
\end{quote}

\subsection*{Your turn}
\begin{itemize}
\raggedright
  \item \emph{Find:} \ta{ஈ} at the front of \textbf{\ta{ஈரம்}}
  \item \emph{Trace it:} \ta{ஈ} three times, saying \emph{ī} as you finish each one
  \item \emph{Write it:} \ta{இ}, then \ta{ஈ}, from memory, naming each
  \item \emph{Write it:} \textbf{\ta{ஈரம்}}
  \item \emph{Read it:} \textbf{\ta{ஈரம்}} aloud once more
\end{itemize}

\begin{tcolorbox}[breakable,colback=teal!4,colframe=teal!35!black,title={Before you move on}]
Which character is this — \ta{ஈ}? (\emph{\textbf{ī}}: the long partner of \textbf{\ta{இ}}.) Which word did it come from? (\textbf{\ta{ஈரம்}}, \emph{īram}, \textquotedblleft{}dampness\textquotedblright{}.) Does it look like \textbf{\ta{இ}}? (\textbf{No}: a bar, two legs, two dots.)
\end{tcolorbox}
\section[\texorpdfstring{\ta{அதன்} \ta{பிறகு} (ataṉ piṟaku)}{அதன் பிறகு (ataṉ piṟaku)}]{\ta{அதன்} \ta{பிறகு} (ataṉ piṟaku) — first, after that, in the end}

\label{lesson:TA-C285-atan-piraku}



\begin{quote}
Say \textquotedblleft{}afterwards\textquotedblright{} (\emph{piṟaku}), \textquotedblleft{}first of all\textquotedblright{} (\emph{mutalil}) and \textquotedblleft{}finally\textquotedblright{} (\emph{kaṭaiciyāka}). Then the word for dampness. (\emph{īram}.) Which letter opens it? (\textbf{\ta{ஈ}}, \emph{ī}.)
\end{quote}

\subsection*{You'll want to know: \ta{முதலில்} … \ta{அதன்} \ta{பிறகு} … \ta{கடைசியில்}}
\begin{quote}
\textbf{\ta{முதலில்} \ta{நான்} \ta{தேநீர்} \ta{குடித்தேன்}. \ta{அதன்} \ta{பிறகு} \ta{பள்ளிக்குப்} \ta{போனேன்}.}
\end{quote}

\begin{quote}
\emph{mutalil nāṉ tēnīr kuṭittēṉ. ataṉ piṟaku paḷḷikkup pōṉēṉ.} — \textquotedblleft{}First I drank tea. After that I went to school.\textquotedblright{}
\end{quote}

\textbf{\ta{முதலில்}} opens a story. \textbf{\ta{அதன்}} is \textbf{\ta{அது}} in the shape that carries endings, \textquotedblleft{}of that\textquotedblright{}, so \textbf{\ta{அதன்} \ta{பிறகு}} is \textquotedblleft{}after that\textquotedblright{}. \textbf{\ta{பிறகு}} alone does the same job.

\begin{quote}
\textbf{\ta{கடைசியில்} \ta{நான்} \ta{இரவு} \ta{பத்து} \ta{மணிக்கு} \ta{வீட்டுக்கு} \ta{வந்தேன்}.}
\end{quote}

\begin{quote}
\emph{kaṭaiciyil nāṉ iravu pattu maṇikku vīṭṭukku vantēṉ.} — \textquotedblleft{}In the end I came home at ten at night.\textquotedblright{}
\end{quote}

\textbf{\ta{கடைசியில்}} (\emph{kaṭaiciyil}, \textquotedblleft{}in the end\textquotedblright{}) is the end of a story, where things came to rest. \textbf{\ta{கடைசியாக}}, \textquotedblleft{}finally, last of all\textquotedblright{}, closes a list. A told story often runs on three posts: \textbf{\ta{முதலில்}} … \textbf{\ta{அதன்} \ta{பிறகு}} … \textbf{\ta{கடைசியில்}}.

\subsection*{You'll want to know: before and after a thing}
\begin{quote}
\textbf{\ta{சாப்பாட்டுக்குப்} \ta{பிறகு} \ta{நாங்கள்} \ta{கடைக்குப்} \ta{போனோம்}.}
\end{quote}

\begin{quote}
\emph{cāppāṭṭukkup piṟaku nāṅkaḷ kaṭaikkup pōṉōm.} — \textquotedblleft{}After the meal we went to the shop.\textquotedblright{}
\end{quote}

\textbf{\ta{சாப்பாடு}} (\emph{cāppāṭu}, \textquotedblleft{}a meal\textquotedblright{}), the noun beside \textbf{\ta{சாப்பிடு}}, takes \textbf{-\ta{க்கு}} before \textbf{\ta{பிறகு}}: \textbf{\ta{சாப்பாட்டுக்குப்} \ta{பிறகு}}, \textquotedblleft{}after the meal\textquotedblright{}. \textbf{\ta{முன்பு}}, \textquotedblleft{}before\textquotedblright{}, takes the same ending: \textbf{\ta{சாப்பாட்டுக்கு} \ta{முன்பு}} (\emph{cāppāṭṭukku muṉpu}, \textquotedblleft{}before the meal\textquotedblright{}).

With a span of days, Tamil counts with \textbf{\ta{கழித்து}}, \textquotedblleft{}having spent\textquotedblright{}, from \textbf{\ta{கழி}}:

\begin{quote}
\textbf{\ta{இரண்டு} \ta{நாள்} \ta{கழித்து} \ta{அவன்} \ta{வந்தான்}.}
\end{quote}

\begin{quote}
\emph{iraṇṭu nāḷ kaḻittu avaṉ vantāṉ.} — \textquotedblleft{}Two days later he came.\textquotedblright{}
\end{quote}

\subsection*{Your turn}
Say these aloud:

\begin{itemize}
\raggedright
  \item \textquotedblleft{}First I drank tea. After that I went to school.\textquotedblright{} — \emph{mutalil nāṉ tēnīr kuṭittēṉ. ataṉ piṟaku paḷḷikkup pōṉēṉ.}
  \item \textquotedblleft{}After the meal we went to the shop\textquotedblright{} — \emph{cāppāṭṭukkup piṟaku nāṅkaḷ kaṭaikkup pōṉōm}
  \item \textquotedblleft{}Two days later he came\textquotedblright{} — \emph{iraṇṭu nāḷ kaḻittu avaṉ vantāṉ}
  \item yesterday in three steps: first, after that, in the end
\end{itemize}

\begin{tcolorbox}[breakable,colback=teal!4,colframe=teal!35!black,title={Before you move on}]
\textquotedblleft{}First\textquotedblright{}? (\textbf{\ta{முதலில்}}.) \textquotedblleft{}After that\textquotedblright{}? (\textbf{\ta{அதன்} \ta{பிறகு}}.) \textquotedblleft{}In the end\textquotedblright{}? (\textbf{\ta{கடைசியில்}}.) \textquotedblleft{}Before the meal\textquotedblright{}? (\textbf{\ta{சாப்பாட்டுக்கு} \ta{முன்பு}}.) \textquotedblleft{}Two days later\textquotedblright{}? (\textbf{\ta{இரண்டு} \ta{நாள்} \ta{கழித்து}}.)
\end{tcolorbox}
\section[\texorpdfstring{\ta{வந்து} (vantu)}{வந்து (vantu)}]{\ta{வந்து} (vantu) — having come, and then}

\label{lesson:TA-C285-vantu}



\begin{quote}
Say \textquotedblleft{}First I drank tea. After that I went to school.\textquotedblright{} (\emph{mutalil nāṉ tēnīr kuṭittēṉ. ataṉ piṟaku paḷḷikkup pōṉēṉ.})

Two sentences, two full verbs. Tamil can fold them into one.
\end{quote}

\subsection*{You'll want to know: the -\ta{உ} and -\ta{இ} form}
\begin{quote}
\textbf{\ta{நான்} \ta{வீட்டுக்கு} \ta{வந்து} \ta{சாப்பிட்டேன்}.}
\end{quote}

\begin{quote}
\emph{nāṉ vīṭṭukku vantu cāppiṭṭēṉ.} — \textquotedblleft{}I came home and ate.\textquotedblright{}
\end{quote}

The first verb carries no person and no time: \textbf{\ta{வந்து}}, \textquotedblleft{}having come\textquotedblright{}. Only the last verb says who and when. So a chain of these runs on to one full verb at the end:

\begin{quote}
\textbf{\ta{அவன்} \ta{கடைக்குப்} \ta{போய்} \ta{ஒரு} \ta{புத்தகம்} \ta{வாங்கி} \ta{வீட்டுக்கு} \ta{வந்தான்}.}
\end{quote}

\begin{quote}
\emph{avaṉ kaṭaikkup pōy oru puttakam vāṅki vīṭṭukku vantāṉ.} — \textquotedblleft{}He went to the shop, bought a book and came home.\textquotedblright{}
\end{quote}

Build it from the \textquotedblleft{}I\textquotedblright{} past. Where the past has \textbf{-\ta{இன்}-} (\textbf{\ta{பேசினேன்}}, \textbf{\ta{வாங்கினேன்}}), the form ends in \textbf{-\ta{இ}}: \textbf{\ta{பேசி}}, \textbf{\ta{வாங்கி}}. Elsewhere, put \textbf{-\ta{உ}} in place of \textbf{-\ta{ஏன்}}: \textbf{\ta{வந்தேன்}} $\to$ \textbf{\ta{வந்து}}, \textbf{\ta{பார்த்தேன்}} $\to$ \textbf{\ta{பார்த்து}}, \textbf{\ta{செய்தேன்}} $\to$ \textbf{\ta{செய்து}}. Two go their own way: \textbf{\ta{போ}} $\to$ \textbf{\ta{போய்}}, \textbf{\ta{சொல்}} $\to$ \textbf{\ta{சொல்லி}}.

You have said it for a long time: \textbf{\ta{போய்} \ta{வருகிறேன்}}, \textquotedblleft{}having gone, I come (back)\textquotedblright{}; \textbf{\ta{தயவுசெய்து}}, \textquotedblleft{}having done (a kindness)\textquotedblright{}.

\subsection*{You'll want to know: -\ta{த} \ta{பிறகு}, after doing}
\begin{quote}
\textbf{\ta{சாப்பிட்ட} \ta{பிறகு} \ta{நான்} \ta{தூங்கினேன்}.}
\end{quote}

\begin{quote}
\emph{cāppiṭṭa piṟaku nāṉ tūṅkiṉēṉ.} — \textquotedblleft{}After eating, I slept.\textquotedblright{}
\end{quote}

\textbf{\ta{சாப்பிட்ட}} is the piece you met in \textbf{\ta{வந்த} \ta{போது}}. Before \textbf{\ta{பிறகு}} it gives \textquotedblleft{}after doing\textquotedblright{}: \textbf{\ta{வந்த} \ta{பிறகு}}, \textquotedblleft{}after coming\textquotedblright{}. The \textbf{-\ta{உ}} / \textbf{-\ta{இ}} form ties two actions tightly; \textbf{-\ta{த} \ta{பிறகு}} marks a clear step in time.

\subsection*{Your turn}
Say these aloud:

\begin{itemize}
\raggedright
  \item \textquotedblleft{}I came home and ate\textquotedblright{} — \emph{nāṉ vīṭṭukku vantu cāppiṭṭēṉ}
  \item \textquotedblleft{}He went to the shop, bought a book and came home\textquotedblright{} — \emph{avaṉ kaṭaikkup pōy oru puttakam vāṅki vīṭṭukku vantāṉ}
  \item \textquotedblleft{}After eating, I slept\textquotedblright{} — \emph{cāppiṭṭa piṟaku nāṉ tūṅkiṉēṉ}
\end{itemize}

\begin{tcolorbox}[breakable,colback=teal!4,colframe=teal!35!black,title={Before you move on}]
\textquotedblleft{}Having come\textquotedblright{}? (\textbf{\ta{வந்து}}.) \textquotedblleft{}Having bought\textquotedblright{}? (\textbf{\ta{வாங்கி}}.) \textquotedblleft{}Having gone\textquotedblright{}? (\textbf{\ta{போய்}}.) Which verb in the chain says who and when? (\textbf{Only the last.}) \textquotedblleft{}After seeing\textquotedblright{}? (\textbf{\ta{பார்த்த} \ta{பிறகு}}.)
\end{tcolorbox}
\section[(reading)]{(reading) — a rainy day}

\label{lesson:TA-C285-mazhai-naal}



\begin{quote}
About eighty words, nearly all met before; the few new forms are glossed after it. The writer is a young man, telling a friend about yesterday.
\end{quote}

\begin{tcolorbox}[breakable,skin=enhanced,colback=orange!6,colframe=orange!45!black,boxrule=0.5pt,arc=1mm,left=6pt,right=6pt,top=4pt,bottom=4pt,fonttitle=\bfseries,title={Read this}]
\begin{quote}
\ta{நேற்று} \ta{நான்} \ta{பாட்டி} \ta{வீட்டுக்குப்} \ta{போனேன்}. \ta{காலையில்} \ta{வெயில்} \ta{இருந்தது}. \ta{நான்} \ta{பேருந்தில்} \ta{போய்} \ta{பாட்டியைப்} \ta{பார்த்தேன்}. \ta{பாட்டி} \ta{எனக்குத்} \ta{தேநீர்} \ta{கொடுத்தார்}. \ta{அவர்} \ta{ஒரு} \ta{காலத்தில்} \ta{ஆசிரியராக} \ta{இருந்தார்}. \ta{அவர்} \ta{எனக்கு} \ta{ஒரு} \ta{கதை} \ta{சொன்னார்}.
\end{quote}

\begin{quote}
\ta{நாங்கள்} \ta{பேசிக்கொண்டிருந்த} \ta{போது} \ta{திடீரென்று} \ta{மழை} \ta{பெய்தது}. \ta{மாலையில்} \ta{நான்} \ta{வீட்டுக்குப்} \ta{புறப்பட்டேன்}. \ta{முதலில்} \ta{பேருந்து} \ta{நிலையத்துக்கு} \ta{நடந்து} \ta{போனேன்}. \ta{என்} \ta{குடை} \ta{வீட்டில்} \ta{இருந்தது}. \ta{அதனால்} \ta{என்} \ta{சட்டை} \ta{ஈரமாக} \ta{இருந்தது}. \ta{என்} \ta{பையிலும்} \ta{ஈரம்}.
\end{quote}

\begin{quote}
\ta{அதன்} \ta{பிறகு} \ta{பேருந்து} \ta{வந்தது}. \ta{பேருந்தில்} \ta{நிறைய} \ta{பயணிகள்} \ta{இருந்தார்கள்}. \ta{கடைசியில்} \ta{இரவு} \ta{எட்டு} \ta{மணிக்கு} \ta{வீட்டுக்கு} \ta{வந்தேன்}. \ta{அப்பா} \ta{செய்தி} \ta{படித்துக்கொண்டிருந்தார்}. \ta{அம்மா} \ta{சமைத்துக்கொண்டிருந்தார்}. \ta{நான்} \ta{சாப்பிட்ட} \ta{பிறகு} \ta{தூங்கினேன்}. \ta{அது} \ta{ஒரு} \ta{நல்ல} \ta{நாள்}.
\end{quote}

Read it once without stopping, then once more, sorting every verb: event, scene, or a step in a chain.
\end{tcolorbox}

\begin{culture}
New forms: \textbf{\ta{பேருந்தில்}}, \textquotedblleft{}by bus\textquotedblright{}; \textbf{\ta{பாட்டியை}}, \textquotedblleft{}grandmother\textquotedblright{}, with the \textbf{-\ta{ஐ}} of the one you see; \textbf{\ta{ஆசிரியராக}}, \textquotedblleft{}(as) a teacher\textquotedblright{}; \textbf{\ta{சொன்னார்}}, \textquotedblleft{}(she) told\textquotedblright{}, the past of \textbf{\ta{சொல்}}; \textbf{\ta{பயணிகள்}}, \textquotedblleft{}passengers\textquotedblright{}; \textbf{\ta{பையிலும்}}, \textquotedblleft{}in the bag too\textquotedblright{}.

The scenes are \textbf{\ta{பேசிக்கொண்டிருந்த} \ta{போது}} and the two \textbf{-\ta{கொண்டிருந்தார்}}; the chains are \textbf{\ta{போய்}} and \textbf{\ta{நடந்து} \ta{போனேன்}}, \textquotedblleft{}went walking\textquotedblright{}. The rest are events.
\end{culture}

\subsection*{Your turn}
\begin{itemize}
\raggedright
  \item \emph{Read it:} the story aloud, once, without stopping
  \item \emph{Say it:} why the shirt was wet, in Tamil — the umbrella was at home
\end{itemize}

\begin{tcolorbox}[breakable,colback=teal!4,colframe=teal!35!black,title={Before you move on}]
Where did he go first in the evening? (\textbf{To the bus station, on foot.}) Why was his shirt wet? (\textbf{His umbrella was at home.}) When did he get home? (\textbf{At eight at night.})
\end{tcolorbox}
\section[(review)]{Review — a story in order}

\label{lesson:TA-R285-a-story-in-order}



\begin{quote}
Say \textquotedblleft{}dampness\textquotedblright{}. (\emph{īram}.) Which letter opens it, and whose long partner is it? (\textbf{\ta{ஈ}}: the long partner of \textbf{\ta{இ}}.) Then \textquotedblleft{}It was raining.\textquotedblright{}

\begin{itemize}
\raggedright
  \item \emph{Write it:} \textbf{\ta{ஈ}} and \textbf{\ta{ஈரம்}}, once each from memory
\end{itemize}
\end{quote}

\subsection*{Your turn}
Say these aloud:

\begin{itemize}
\raggedright
  \item First I went to Grandmother's house. After that we drank tea.
  \item While we were talking, suddenly it rained.
  \item I walked to the bus station. My shirt was wet.
  \item In the end I came home and ate.
  \item a rainy day of your own, in four sentences, in order
\end{itemize}

\begin{tcolorbox}[breakable,colback=teal!4,colframe=teal!35!black,title={Before you move on}]
\textquotedblleft{}Dampness\textquotedblright{}? (\textbf{\ta{ஈரம்}}.) \textquotedblleft{}First … after that … in the end\textquotedblright{}? (\textbf{\ta{முதலில்}} … \textbf{\ta{அதன்} \ta{பிறகு}} … \textbf{\ta{கடைசியில்}}.) \textquotedblleft{}Having come\textquotedblright{}? (\textbf{\ta{வந்து}}.) \textquotedblleft{}After eating\textquotedblright{}? (\textbf{\ta{சாப்பிட்ட} \ta{பிறகு}}.)
\end{tcolorbox}
